% Options for packages loaded elsewhere
\PassOptionsToPackage{unicode}{hyperref}
\PassOptionsToPackage{hyphens}{url}
\PassOptionsToPackage{dvipsnames,svgnames,x11names}{xcolor}
%
\documentclass[
  11pt,
]{article}
\usepackage{amsmath,amssymb}
\usepackage{iftex}
\ifPDFTeX
  \usepackage[T1]{fontenc}
  \usepackage[utf8]{inputenc}
  \usepackage{textcomp} % provide euro and other symbols
\else % if luatex or xetex
  \usepackage{unicode-math} % this also loads fontspec
  \defaultfontfeatures{Scale=MatchLowercase}
  \defaultfontfeatures[\rmfamily]{Ligatures=TeX,Scale=1}
\fi
\usepackage{lmodern}
\ifPDFTeX\else
  % xetex/luatex font selection
\fi
% Use upquote if available, for straight quotes in verbatim environments
\IfFileExists{upquote.sty}{\usepackage{upquote}}{}
\IfFileExists{microtype.sty}{% use microtype if available
  \usepackage[]{microtype}
  \UseMicrotypeSet[protrusion]{basicmath} % disable protrusion for tt fonts
}{}
\makeatletter
\@ifundefined{KOMAClassName}{% if non-KOMA class
  \IfFileExists{parskip.sty}{%
    \usepackage{parskip}
  }{% else
    \setlength{\parindent}{0pt}
    \setlength{\parskip}{6pt plus 2pt minus 1pt}}
}{% if KOMA class
  \KOMAoptions{parskip=half}}
\makeatother
\usepackage{xcolor}
\usepackage[margin=1in]{geometry}
\usepackage{longtable,booktabs,array}
\usepackage{calc} % for calculating minipage widths
% Correct order of tables after \paragraph or \subparagraph
\usepackage{etoolbox}
\makeatletter
\patchcmd\longtable{\par}{\if@noskipsec\mbox{}\fi\par}{}{}
\makeatother
% Allow footnotes in longtable head/foot
\IfFileExists{footnotehyper.sty}{\usepackage{footnotehyper}}{\usepackage{footnote}}
\makesavenoteenv{longtable}
\setlength{\emergencystretch}{3em} % prevent overfull lines
\providecommand{\tightlist}{%
  \setlength{\itemsep}{0pt}\setlength{\parskip}{0pt}}
\setcounter{secnumdepth}{-\maxdimen} % remove section numbering
\ifLuaTeX
  \usepackage{selnolig}  % disable illegal ligatures
\fi
\IfFileExists{bookmark.sty}{\usepackage{bookmark}}{\usepackage{hyperref}}
\IfFileExists{xurl.sty}{\usepackage{xurl}}{} % add URL line breaks if available
\urlstyle{same}
\hypersetup{
  pdfauthor={Stefano Alimonti --- stefalimonti@gmail.com},
  colorlinks=true,
  linkcolor={blue},
  filecolor={Maroon},
  citecolor={Blue},
  urlcolor={blue},
  pdfcreator={LaTeX via pandoc}}

\author{Stefano Alimonti --- stefalimonti@gmail.com}
\date{March 2026}

\begin{document}

{
\hypersetup{linkcolor=black}
\setcounter{tocdepth}{2}
\tableofcontents
}
\section{The Angular Uniqueness of Base 2 in Positional
Multiplication}\label{the-angular-uniqueness-of-base-2-in-positional-multiplication}

\textbf{Author:} Stefano Alimonti \textbf{Affiliation:} Independent Researcher
\textbf{Date:} March 2026

\begin{center}\rule{0.5\linewidth}{0.5pt}\end{center}

\newpage

\subsection{Abstract}\label{abstract}

We prove that base 2 is the unique integer base for which the D-parity boundary
in positional multiplication becomes a straight line in angular coordinates:
\((1+\tan\alpha)(1+\tan\beta) = b\) reduces to \(\alpha + \beta = \pi/4\) if and
only if \(b = 2\). This geometric fact explains why \(\pi\) enters the
conjectured trace anomaly \(c_1 = \pi/18\) for binary multiplication --- the
integration domain is a right triangle with a \(\pi/4\) hypotenuse --- and
predicts that \(c_1(b)\) for \(b > 2\) does not involve \(\pi\).

\textbf{Keywords:} angular coordinates, D-parity boundary, trace anomaly,
carries, positional multiplication, base dependence

\textbf{MSC 2020:} 11A63, 60C05, 26B15

\begin{center}\rule{0.5\linewidth}{0.5pt}\end{center}

\newpage

\subsection{1. Introduction}\label{introduction}

The trace anomaly

\[c_1 := \mathbb{E}[c_{\mathrm{top}-1}] - 1\]

governs the leading correction in the carry--zeta approximation {[}B{]}:

\[\ln R(l,s) = \frac{c_1}{l^s} + O(l^{-2\sigma}).\]

For base 2, high-statistics Monte Carlo and exact enumeration up to \(K = 21\)
establish \(c_1 = \pi/18 = 0.17453\ldots\) to \({\sim}4.3\) digits from direct
enumeration {[}E{]}. The appearance of \(\pi\) in a problem of positional
arithmetic is surprising --- no trigonometric structure is apparent in the carry
recurrence

\[c_{k+1} = \lfloor(\mathrm{conv}_k + c_k)/2\rfloor\]

. (The trace anomaly \(c_1\) of the Euler product correction studied in this
paper and in {[}B{]} is the leading spectral constant of the carry correction
factor; it is related to, but distinct from, the macroscopic cascade trace
anomaly \(\Delta R\) studied in {[}E{]}, which concerns the full sector ratio of
the cascade valuation.)

This paper explains where \(\pi\) comes from. We introduce angular coordinates
\((\alpha, \beta)\) for the normalized factors and show that the D-parity
boundary --- the hyperbola \((1+X)(1+Y) = b\) separating even-length from
odd-length products --- becomes a straight line \(\alpha + \beta = \pi/4\) in
these coordinates if and only if \(b = 2\). The constant \(\pi/4\) in the
integration limit pulls \(\pi\) into the trace anomaly through the geometry of
the domain, not through the arithmetic of carries.

For \(b > 2\) the boundary is a curve, the angular range involves
\(\arctan(b-1)\) (irrational multiples of \(\pi\)), and multi-base experiments
confirm that \(c_1(b)\) departs sharply from \(\pi/18\).

\begin{center}\rule{0.5\linewidth}{0.5pt}\end{center}

\newpage

\subsection{2. Angular Coordinates and the D-Parity
Boundary}\label{angular-coordinates-and-the-d-parity-boundary}

\subsubsection{2.1 Setup}\label{setup}

Let \(p, q\) be independent random \(d\)-digit integers in base \(b\), with
leading digit \(\geq 1\) and remaining digits uniform on
\(\lbrace{}0, \ldots, b-1\rbrace{}\). Write \(p = b^{d-1}(1+X)\),
\(q = b^{d-1}(1+Y)\), so that \(X, Y \in [0, b-1)\) are continuous
approximations to the fractional parts. The product \(N = pq\) has

\[N = b^{2d-2}\,W, \qquad W = (1+X)(1+Y) \in [1, b^2).\]

The digit count \(D = \lfloor\log_b N\rfloor + 1\) satisfies:

\begin{itemize}
\tightlist
\item
  \(D = 2d\) (even) when \(W \geq b\),
\item
  \(D = 2d - 1\) (odd) when \(W < b\).
\end{itemize}

This is the \textbf{D-parity} decomposition. The boundary is the hyperbola
\((1+X)(1+Y) = b\) in the unit square \([0, b-1)^2\).

\subsubsection{2.2 Angular Substitution}\label{angular-substitution}

Introduce angular coordinates via \(X = \tan\alpha\), \(Y = \tan\beta\), where
\(\alpha, \beta \in [0, \arctan(b-1))\). The Jacobian is
\(dX\,dY = \sec^2\!\alpha\;\sec^2\!\beta\;d\alpha\,d\beta\), and

\[W = (1 + \tan\alpha)(1 + \tan\beta).\]

\textbf{Proposition (Angular boundary).} The D-parity boundary
\((1+\tan\alpha)(1+\tan\beta) = b\) in angular coordinates satisfies

\[\tan(\alpha + \beta) = \frac{b - 1 - P}{1 - P}, \qquad P = \tan\alpha\,\tan\beta.\]

\emph{Proof.} Expand the left side:
\(1 + \tan\alpha + \tan\beta + \tan\alpha\,\tan\beta = b\), so
\(\tan\alpha + \tan\beta = b - 1 - P\). The addition formula gives
\(\tan(\alpha+\beta) = (\tan\alpha + \tan\beta)/(1 - \tan\alpha\,\tan\beta) = (b-1-P)/(1-P)\).
\(\square\)

\subsubsection{2.3 The Angular Uniqueness
Theorem}\label{the-angular-uniqueness-theorem}

\textbf{Theorem (Angular Uniqueness).} The expression \((b-1-P)/(1-P)\) is
independent of \(P\) if and only if \(b = 2\). Consequently:

\begin{itemize}
\tightlist
\item
  For \(b = 2\): the D-parity boundary is the straight line
  \(\alpha + \beta = \pi/4\).
\item
  For \(b > 2\): the boundary is a curve, and the angular sum \(\alpha + \beta\)
  varies along it.
\end{itemize}

\emph{Proof.} Write \(f(P) = (b-1-P)/(1-P)\). Compute the derivative:

\[f'(P) = \frac{-(1-P) + (b-1-P)}{(1-P)^2} = \frac{b - 2}{(1-P)^2}.\]

The numerator \(b - 2\) vanishes if and only if \(b = 2\). When \(b = 2\):
\(f(P) = (1-P)/(1-P) = 1\) for all \(P \in [0,1)\), so
\(\tan(\alpha+\beta) = 1\), giving \(\alpha + \beta = \pi/4\). When \(b > 2\):
\(f'(P) > 0\), so \(\tan(\alpha+\beta)\) increases with \(P\) and the boundary
is not a level set of \(\alpha + \beta\). \(\square\)

\textbf{Remark.} The algebraic mechanism is transparent: \(b - 1 = 1\) is the
unique case where the numerator and denominator of \((b-1-P)/(1-P)\) are
proportional, producing a cancellation that eliminates \(P\).

\begin{center}\rule{0.5\linewidth}{0.5pt}\end{center}

\newpage

\subsection{3. Consequences}\label{consequences}

\subsubsection{3.1 The D-Odd Region for Base
2}\label{the-d-odd-region-for-base-2}

For \(b = 2\), the D-odd condition \(W < 2\) maps to the right triangle

\[\mathcal{T} = \lbrace{}(\alpha, \beta) : \alpha \geq 0,\; \beta \geq 0,\; \alpha + \beta < \pi/4\rbrace{}\]

in the angular square \([0, \pi/4)^2\). The D-odd probability is

\[P_o = \int_0^{\pi/4}\!\int_0^{\pi/4 - \alpha} \sec^2\!\alpha\;\sec^2\!\beta\;d\beta\;d\alpha.\]

Evaluating the inner integral:

\[\int_0^{\pi/4 - \alpha}\sec^2\!\beta\;d\beta = \tan(\pi/4 - \alpha) = \frac{1 - \tan\alpha}{1 + \tan\alpha}.\]

Substituting \(t = \tan\alpha\):

\[P_o = \int_0^1 \frac{1-t}{1+t}\,dt = \bigl[2\ln(1+t) - t\bigr]_0^1 = 2\ln 2 - 1 \approx 0.3863.\]

The integration limit \(\pi/4\) is the source of \(\pi\) in all subsequent
integrals over the D-odd domain.

\subsubsection{3.2 The Inner Integral on
Isochrones}\label{the-inner-integral-on-isochrones}

Introduce the isochrone coordinate \(s = \alpha + \beta\). Along the level set
at angular distance \(s\) from the vertex, the density integral is

\[I(s) = \int_0^{s} \sec^2\!\alpha\;\sec^2(s - \alpha)\;d\alpha = \frac{2(1+T^2)\bigl[T^2 - \ln(1+T^2)\bigr]}{T^3}, \qquad T = \tan s.\]

\emph{Derivation.} Write
\(\sec^2\!\alpha\;\sec^2(s-\alpha) = (1 + \tan^2\!\alpha)(1 + \tan^2(s-\alpha))\).
Substituting \(u = \tan\alpha\) with \(du = \sec^2\!\alpha\;d\alpha\) and using
\(\tan(s-\alpha) = (T - u)/(1 + Tu)\):

\[I(s) = \int_0^{T} \left(1 + \frac{(T-u)^2}{(1+Tu)^2}\right) du = \int_0^T \frac{(1+Tu)^2 + (T-u)^2}{(1+Tu)^2}\,du.\]

Expanding the numerator: \((1+Tu)^2 + (T-u)^2 = (1 + T^2)(1 + u^2)\), so

\[I(s) = (1+T^2)\int_0^{T} \frac{1 + u^2}{(1+Tu)^2}\,du.\]

Substituting \(w = 1 + Tu\) and separating partial fractions:

\[\int_0^T \frac{1+u^2}{(1+Tu)^2}\,du = \frac{1}{T^3}\Bigl[-(T^2{+}1)/w - 2\ln w + w\Bigr]_1^{1+T^2} = \frac{2\bigl[T^2 - \ln(1+T^2)\bigr]}{T^3},\]

yielding the stated closed form.

\textbf{Verification.} \(P_o = \int_0^{\pi/4} I(s)\,ds = 2\ln 2 - 1\), confirmed
analytically by reversing the order of integration.

\subsubsection{\texorpdfstring{3.3 Representation of
\(c_1\)}{3.3 Representation of c\_1}}\label{representation-of-c_1}

The trace anomaly admits a single-integral representation over the angular
domain:

\[c_1 = \int_0^{\pi/4} \langle c_{M-1} - 1\rangle(s) \cdot I(s)\;ds + (1 + 3\ln(3/4))\]

where \(\langle c_{M-1} - 1\rangle(s)\) is the isochrone average of
\(c_{M-1} - 1\) at angular distance \(s\), and the constant \(1 + 3\ln(3/4)\) is
the D-even contribution (the D-even cascade is degenerate {[}E, §8.1{]}, making
the schoolbook computation exact and the D-even \(c_1\) component analytically
closed {[}B{]}). The upper limit \(\pi/4\) --- a direct consequence of the
Angular Uniqueness Theorem --- is the sole entry point for \(\pi\).

\subsubsection{\texorpdfstring{3.4 Base \(b > 2\): Curved
Boundaries}{3.4 Base b \textgreater{} 2: Curved Boundaries}}\label{base-b-2-curved-boundaries}

For \(b > 2\), the angular sum along the D-parity boundary varies from
\(\arctan(b-1)\) (at \(P = 0\), i.e., \(\alpha = 0\) or \(\beta = 0\)) upward.
The value \(\arctan(b-1)\) is not a rational multiple of \(\pi\) for any integer
\(b > 2\) (by the Niven--Mann theorem: \(\arctan(n)\) is a rational multiple of
\(\pi\) only for \(n = 0, \pm 1\)). The integration domain is bounded by a
curve, not a line, and the integration \emph{limits} no longer produce clean
factors of \(\pi\).

\textbf{Conjecture.} \(c_1(b)\) for \(b > 2\) does not involve \(\pi\).

\emph{Supporting evidence:} (i) Niven--Mann constrains the integration limits;
(ii) base-3 numerical data (\(c_1(3) \approx 0.5985\), candidate
\(\ln 3 - 1/2\)) shows no \(\pi\) involvement (§4.1); (iii) PSLQ analysis on
base-3 data excludes \(\pi\) from the constant basis at \({\sim}4\)-digit
precision (G03, G12). \emph{Caveat:} this argument constrains the limits but not
the integrand --- in principle, \(\pi\) could enter through the carry kernel
\(\langle c_{M-1} - 1 \rangle(s)\) or through combinations of arctan values. A
proof would require showing the full integral, not just the limits, is
\(\pi\)-free.

\begin{center}\rule{0.5\linewidth}{0.5pt}\end{center}

\newpage

\subsection{4. Experimental Verification}\label{experimental-verification}

\subsubsection{4.1 Multi-Base Monte Carlo}\label{multi-base-monte-carlo}

Monte Carlo simulation with \(10^8\) samples per base (\(d = 64\)-bit factors)
yields:

\begin{longtable}[]{@{}llll@{}}
\toprule\noalign{}
Base \(b\) & \(c_1(b)\) & Std. error & Consistent with \(\pi/18\)? \\
\midrule\noalign{}
\endhead
\bottomrule\noalign{}
\endlastfoot
2 & 0.17465 & \(4.4 \times 10^{-5}\) & Yes (\(0.3\sigma\)) \\
3 & 0.599 & \(1 \times 10^{-3}\) & No (\(>400\sigma\)) \\
5 & 1.61 & \(2 \times 10^{-3}\) & No (\(>700\sigma\)) \\
7 & 2.64 & \(3 \times 10^{-3}\) & No (\(>800\sigma\)) \\
10 & 4.20 & \(4 \times 10^{-3}\) & No (\(>1000\sigma\)) \\
\end{longtable}

The trace anomaly grows with base and departs decisively from \(\pi/18\) for all
\(b > 2\).

\subsubsection{4.2 Exact Enumeration for Base
3}\label{exact-enumeration-for-base-3}

Exact enumeration of \(c_1(3, K)\) for \(K = 2, \ldots, 12\)
(\(3^{24} \approx 2.82 \times 10^{11}\) configurations at \(K = 12\),
parallelized with OpenMP) yields:

\begin{longtable}[]{@{}llll@{}}
\toprule\noalign{}
\(K\) & \(c_1(3, K)\) & \(\Delta\) & \(\rho(K)\) \\
\midrule\noalign{}
\endhead
\bottomrule\noalign{}
\endlastfoot
2 & 0.50000 & --- & --- \\
3 & 0.48148 & \(-1.85 \times 10^{-2}\) & --- \\
4 & 0.53185 & \(+5.04 \times 10^{-2}\) & 2.72 \\
5 & 0.54670 & \(+1.48 \times 10^{-2}\) & 0.295 \\
6 & 0.56709 & \(+2.04 \times 10^{-2}\) & 1.37 \\
7 & 0.57461 & \(+7.52 \times 10^{-3}\) & 0.369 \\
8 & 0.58280 & \(+8.19 \times 10^{-3}\) & 1.09 \\
9 & 0.58684 & \(+4.04 \times 10^{-3}\) & 0.493 \\
10 & 0.59015 & \(+3.31 \times 10^{-3}\) & 0.819 \\
11 & 0.59822 & \(+8.07 \times 10^{-3}\) & 2.44 \\
12 & 0.59851 & \(+2.84 \times 10^{-4}\) & 0.035 \\
\end{longtable}

The \(K = 11, 12\) values are exact rationals:
\(c_1(3, 11) = 8342998915 / 13946313395\) and
\(c_1(3, 12) = 75125607729 / 125521846608\).

A dramatic \textbf{odd-even oscillation} in \(\rho(K)\) dominates the
convergence: odd \(K\) overshoots, even \(K\) undershoots, with the oscillation
amplitude decreasing. This makes multi-term extrapolation unstable. The 1-term
Richardson extrapolation from the most recent pair gives
\(c_\infty^{\text{Rich}}(11, 12) = 0.59865\), close to
\(\ln(3) - 1/2 = 0.59861\ldots\) (\(\Delta = 3.6 \times 10^{-5}\),
\(\sim 1.4\sigma\)).

PSLQ search with extended basis
\(\lbrace{}1, \ln 2, \ln 3, \ln^2 3, \arctan(1/2), \zeta(2)/9, \sqrt{3}, \pi\rbrace{}\)
finds no exact relation at the available precision (\(\sim 5\)--\(6\) digits).
The candidate \(\ln(3) - 1/2\) is suggestive but not confirmed.

\textbf{Enhanced analysis (G12).} Separating \(K\) into even
\(\lbrace{}2,4,6,8,10,12\rbrace{}\) and odd \(\lbrace{}3,5,7,9,11\rbrace{}\)
subsequences eliminates the oscillation (each subsequence converges
monotonically), consistent with a \((-1/3)^K\) alternating eigenvalue. However,
multi-term Richardson extrapolation on each subsequence with rate \(1/3\)
overshoots the target by \(\sim 10^{-3}\): the effective correction rate within
the even subsequence is \(r_{\text{eff}} \approx 1.07\), not \(1/9\). This
reveals a \textbf{pre-asymptotic structure}: the correction coefficients \(A_K\)
depend on \(K\) (the boundary observable changes with each cascade depth), so
the true expansion is \(c_1(K) = c_\infty + (A_0 + A_1/K + \cdots)(1/3)^K\),
defeating simple geometric Richardson. Nonlinear models (alternating eigenvalues
\(\pm 1/3, \pm 1/9, 1/27\)) also underfit (RMS \(\sim 5 \times 10^{-3}\)). All
estimates cluster in \([0.590, 0.600]\), bracketing \(\ln(3) - 1/2 = 0.5986\)
(consistent at \(0.6\sigma\)), but the available precision (\(\sim 2\)--\(3\)
digits) is far below the \(\sim 8\) digits required for definitive PSLQ. Data at
\(K \geq 16\) would be needed.

\subsubsection{4.3 Convergence Rate}\label{convergence-rate}

For base 2, the effective convergence rate \(\rho(K)\) approaches \(1/2\)
monotonically (\(\rho(17) = 0.505\), proved to equal \(\lambda_2 = 1/b\) via
transfer operator analysis).

For base 3, the convergence shows a pronounced \textbf{odd-even oscillation}:
\(\rho(K)\) alternates between values well above 1 (at odd \(K\)) and well below
1 (at even \(K\)). The sequence \(\rho(K)\) for \(K = 4, \ldots, 12\) is:
\(2.72, 0.30, 1.37, 0.37, 1.09, 0.49, 0.82, 2.44, 0.035\). This oscillation
indicates complex sub-dominant eigenvalues in the transfer operator, consistent
with the \((b-1) = 2\) off-diagonal channels in the base-3 carry Markov chain.

Despite the oscillation, the \textbf{geometric mean} of consecutive \(\rho\)
values decreases toward \(1/3\):

\begin{longtable}[]{@{}ll@{}}
\toprule\noalign{}
Pair & \(\sqrt{\rho_{\text{odd}} \cdot \rho_{\text{even}}}\) \\
\midrule\noalign{}
\endhead
\bottomrule\noalign{}
\endlastfoot
\((K=4,5)\) & 0.90 \\
\((K=6,7)\) & 0.71 \\
\((K=8,9)\) & 0.73 \\
\((K=10,11)\) & 1.41 \\
\((K=11,12)\) & 0.29 \\
\end{longtable}

The last pair (\(K = 11, 12\)) undershoots \(1/3\), suggesting the oscillation
is damping. More data (\(K \geq 14\)) would be needed to confirm the asymptotic
geometric-mean rate equals \(1/3 = 1/b\).

This identifies the Diaconis--Fulman sub-dominant eigenvalue \(\lambda_2 = 1/b\)
as the universal convergence bottleneck, though for \(b > 2\) the approach to
this rate involves oscillatory transients from complex eigenvalues.

\subsubsection{4.4 Angular Boundary Curvature}\label{angular-boundary-curvature}

The standard deviation \(\sigma(\alpha + \beta)\) along the D-parity boundary,
sampled at 1000 uniformly spaced points, quantifies the departure from
linearity:

\begin{longtable}[]{@{}ll@{}}
\toprule\noalign{}
Base \(b\) & \(\sigma(\alpha + \beta)\) \\
\midrule\noalign{}
\endhead
\bottomrule\noalign{}
\endlastfoot
2 & 0.000 (exact) \\
3 & 0.046 \\
5 & 0.137 \\
10 & 0.249 \\
100 & 0.413 \\
\end{longtable}

The curvature vanishes only for \(b = 2\) and grows monotonically, confirming
the theorem.

\subsubsection{4.5 Transfer Operator Spectrum}\label{transfer-operator-spectrum}

The carry transfer operator for base-\(b\) multiplication, acting on carry
states \(\lbrace{}0, \ldots, b-1\rbrace{}\), has eigenvalues
\(\lambda_k = 1/b^k\) for \(k = 0, 1, \ldots, b-1\) --- exactly the
Diaconis--Fulman spectrum. This holds for all tested bases
\(b = 2, 3, 5, 7, 10\) and extends the Diaconis--Fulman theorem from addition
carries to multiplication carries.

The mechanism is the division by \(b\) in the carry recurrence: a linear
perturbation \(\delta(c) = c - \mathbb{E}[c]\) maps to \(\delta/b\) at each
step, regardless of the convolution distribution. Higher-order perturbations
decay at rates \(1/b^2, 1/b^3, \ldots\), completing the spectrum.

\begin{center}\rule{0.5\linewidth}{0.5pt}\end{center}

\newpage

\subsection{5. Discussion}\label{discussion}

\subsubsection{\texorpdfstring{5.1 The Origin of
\(\pi\)}{5.1 The Origin of \textbackslash pi}}\label{the-origin-of-pi}

The Angular Uniqueness Theorem explains why \(\pi\) enters \(c_1(2) = \pi/18\).
The constant \(\pi\) does not come from the carry arithmetic --- neither the
carry recurrence nor the convolution sums involve trigonometric quantities.
Instead, \(\pi\) enters through the geometry of the integration domain: the
D-odd region, when expressed in the natural angular coordinates for the
multiplicative structure \((1+X)(1+Y)\), is a right triangle with hypotenuse
\(\alpha + \beta = \pi/4\). The integration limit \(\pi/4\) propagates into the
final answer.

For \(b > 2\), the boundary is curved, the domain is not a simplex, and the
integration limits involve \(\arctan(b-1)\) --- quantities that are irrational
multiples of \(\pi\). The limits alone do not produce clean factors of \(\pi\);
whether \(\pi\) could enter through the integrand remains an open question (see
Conjecture in §3.4).

\subsubsection{5.2 Geometry Versus Algebra}\label{geometry-versus-algebra}

The trace anomaly \(c_1\) depends on two distinct structures:

\begin{enumerate}
\def\labelenumi{\arabic{enumi}.}
\item
  \textbf{The spectrum} --- the transfer operator eigenvalues
  \(\lambda_k = 1/b^k\) --- which governs the convergence rate
  \(\rho(K) \to 1/b\). This is algebraic, universal across bases, and proved for
  both addition and multiplication carries.
\item
  \textbf{The boundary geometry} --- the shape of the D-parity region in angular
  coordinates --- which determines the value of \(c_1\) itself. This is
  geometric, base-specific, and uniquely simple for \(b = 2\).
\end{enumerate}

The formula \(c_1(K) = \pi/18 + P(K) \cdot (1/2)^K\) from {[}B, experiment
B30{]} separates these contributions cleanly: \(\pi/18\) is geometric (the
integration domain), \(1/2\) is algebraic (the spectral gap), and the polynomial
prefactor \(P(K)\) arises from non-Markov digit-bit correlations in the
multiplication carry chain.

\subsubsection{5.3 Connections}\label{connections}

The Angular Uniqueness Theorem complements the spectral theory of carries
{[}A{]}, which establishes the Diaconis--Fulman eigenvalue structure, and the
carry--zeta approximation framework {[}B{]}, which uses \(c_1\) as the leading
correction constant. The present result explains why the correction involves
\(\pi\) for the base (\(b = 2\)) in which the framework was originally
developed, and predicts that multi-base extensions will encounter different ---
and likely more complex --- transcendental structure.

The trace anomaly analysis {[}E{]} established \(c_1 = \pi/18\) to \({\sim}4.3\)
digits from direct enumeration (\(K = 21\)); \(\Sigma_{\text{even}}\) achieves
\(7.3\) digits via Richardson extrapolation. A formal proof remains the central
open problem. The Angular Uniqueness Theorem identifies the geometric origin of
\(\pi\) and thereby constrains the form any such proof must take: it must
ultimately evaluate an integral over the triangular domain
\(\alpha + \beta < \pi/4\).

\begin{center}\rule{0.5\linewidth}{0.5pt}\end{center}

\newpage

\subsection{6. Reproducibility}\label{reproducibility}

All scripts are in \texttt{experiments/}, following the naming convention
\texttt{G\{NN\}\_\{description\}.\{py\textbar{}c\}}. Core experiments: G01--G12;
extensions G13--G15 are documented in \texttt{experiments/README.md}.

\begin{longtable}[]{@{}
  >{\raggedright\arraybackslash}p{(\columnwidth - 4\tabcolsep) * \real{0.2581}}
  >{\raggedright\arraybackslash}p{(\columnwidth - 4\tabcolsep) * \real{0.3226}}
  >{\raggedright\arraybackslash}p{(\columnwidth - 4\tabcolsep) * \real{0.4194}}@{}}
\toprule\noalign{}
\begin{minipage}[b]{\linewidth}\raggedright
Script
\end{minipage} & \begin{minipage}[b]{\linewidth}\raggedright
Language
\end{minipage} & \begin{minipage}[b]{\linewidth}\raggedright
Description
\end{minipage} \\
\midrule\noalign{}
\endhead
\bottomrule\noalign{}
\endlastfoot
G01 & C & High-precision Monte Carlo for \(c_1(3)\) (\(10^8\)--\(10^9\)
samples) \\
G02 & C & Exact enumeration \(c_1(3, K)\) for \(K = 2\)--\(12\) (OpenMP) \\
G03 & Python & PSLQ analysis on extrapolated \(c_1(3)\) \\
G04 & Python & Transfer operator eigenvalue verification, 5 bases (§4.5) \\
G05--G06 & Python & Base-3 eigenvalue verification and extensions \\
G07 & C & Base-3 exact enumeration \(K = 2\)--\(10\) \\
G08--G09 & Python & Multi-term Diaconis--Fulman fit; PSLQ on \(c_1(3)\) \\
G10--G11 & C/Python & Extended enumeration \(K = 11\)--\(14\); enhanced PSLQ
with 5-term Richardson \\
G12 & Python & Odd-even separation, alternating eigenvalue models \\
\end{longtable}

Key result from G11: base-3 transfer operators confirm the full Diaconis--Fulman
spectrum \(\lbrace{}1, 1/3, 1/9, 1/27\rbrace{}\). Best estimate
\(c_1(3) \approx 0.5985\), candidate \(\ln 3 - 1/2\), PSLQ inconclusive at
4-digit precision. Definitive identification requires \(K \geq 16\).

Requirements: Python 3.8+, NumPy, SciPy, mpmath. C compiler (gcc/clang) for
\texttt{.c} scripts.

\begin{center}\rule{0.5\linewidth}{0.5pt}\end{center}

\newpage

\subsection{7. References}\label{references}

\begin{enumerate}
\def\labelenumi{\arabic{enumi}.}
\tightlist
\item
  P. Diaconis, J. Fulman, ``Carries, Shuffling, and Symmetric Functions,''
  \emph{Adv. Appl. Math.} 43(2), 176--196, 2009.
\item
  I. Niven, ``Irrational Numbers,'' \emph{Carus Mathematical Monographs} 11,
  MAA, 1956.
\item
  {[}A{]} Companion paper: ``Spectral Theory of Carries in Positional
  Multiplication,'' this series. doi:10.5281/zenodo.18895593
\item
  {[}B{]} Companion paper: ``Carry Polynomials and the Euler Product: An
  Approximation Framework,'' this series. doi:10.5281/zenodo.18895597
\item
  {[}P1{]} Companion paper: ``π from Pure Arithmetic: A Spectral Phase
  Transition in the Binary Carry Bridge,'' this series.
  doi:10.5281/zenodo.18895611
\item
  {[}P2{]} Companion paper: ``The Sector Ratio in Binary Multiplication: From
  Markov Failure to Transcendence,'' this series. doi:10.5281/zenodo.18895615
\item
  {[}E{]} Companion paper: ``The Trace Anomaly of Binary Multiplication,'' this
  series. doi:10.5281/zenodo.18895604
\item
  {[}F{]} Companion paper: ``Exact Covariance Structure of Binary Carry
  Chains,'' this series. doi:10.5281/zenodo.18895607
\end{enumerate}

\begin{center}\rule{0.5\linewidth}{0.5pt}\end{center}

\emph{CC BY 4.0. Code: MIT License.}

\end{document}
